\scanpage{51}
\mspage{\unsure{7}}
% archive sheet number: 50
 d'une structure déjà orientée si $a$ est une arête libre (correspondant à un $C\in\pi_0(\partial X)$), ou \unsure{bien} disposer de l'\unsure{ens.|une} des deux orientations $\omega(a)$ -- dans le cas où $a$ est une arête ordinaire, indique que l'on tient en mains, pour chaque $a\in A$, un polygone combinatoire $P_a$, dont l'ens des sommets est $R_a$, et l'ens des arêtes peut se décrire \deleted{comme} de façon canonique en termes de $R_a,\omega_a$ et $\omega_a\longrightarrow\mathrm{Circ}(R_a)$. Il peut être plus commode de la donner d'emblée dans la structure, où la famille des $R_a$ serait remplacée par un \deleted{qui serait donc un certain} \uline{carton combinatoire}, i.e. un ensemble de repères $\underline{\vec R}$ (dits ``repères de rigidification''), avec une opération de
\[
\unsure{\mathbb D|D}_\infty=\{\sigma_0,\sigma_1\mid\sigma_0^2=\sigma_1^2=1\}
\]
dessus, et une application
\diagram{\begin{tikzcd}[ampersand replacement=\&]
\underline{\vec R}\arrow[r] \&\vec A
\end{tikzcd}}
compatible avec l'hom
\diagram{\begin{tikzcd}[ampersand replacement=\&]
\mathbb D_\infty\arrow[r,"\unsure{\epsilon|\varepsilon}"] \&\{\pm1\}
\end{tikzcd}}
sur les groupes d'opérateurs, mais de plus un sous ens
\[
\underline{\vec R}^{+}_\ell\subset\underline{\vec R}\mid\vec A_\ell
\]
\margin{($\vec A_\ell={}$arêtes libres, correspondant aux cercles de recollement, i.e. les $\in\pi_0(\deleted{\unsure{\sigma}}\partial X)$)}
définissant une orientation des cartons correspondants.
